\chapter{Experimental Design and Reproducibility Methodology}

\section{Research questions and controlled axes}

The experiments ask four bounded questions. First, how do stored states change
when arithmetic precision or effective build semantics change? Second, within
matched configurations, how do the HLL/HLLD choice and CPU/GPU execution path
affect output and elapsed time? Third, how does the fp32--fp64 discrepancy vary
with resolution and time, and how does it compare with virtual-precision
stochastic variation? Fourth, do thread count and CFL number alter the covered
observations? Precision, build semantics, solver, device, resolution and time
are primary axes. Thread count and CFL number are robustness axes. MPI,
additional GPU solvers and cross-machine performance are deliberately deferred.

\section{Test-case matrix}

The ideal-MHD matrix contains the one-dimensional Brio--Wu shock tube and the
two-dimensional periodic Orszag--Tang and Kelvin--Helmholtz problems. These
cases cover a discontinuous one-dimensional flow, a two-dimensional interacting
wave system and a shear instability. Sod and Liska--Wendroff Configuration~3
provide a compact one- and two-dimensional Euler continuity check for the
cross-system comparison; they do not repeat the Euler validation study from
Report~1. Each result states its case, grid, final time, solver and baseline,
so observations are not transferred between unmatched configurations.

\section{Build and run matrix}

The arithmetic builds use binary64 or binary32 storage and computation. Matched
build comparisons additionally vary the effective optimisation, floating-point
math and branch-rule semantics recorded by the harness. Directory names are
identifiers rather than evidence of compiler behaviour. The isolated MSVC
matrix compares /O2 with /Ox, compiler-default math with /fp:fast at /O2, and
$\leq$ with $<$ at /O2 and default math, changing one recorded axis per pair.
Deterministic
precision comparisons use the same-solver fp64 run at the same grid and final
time as their numerical baseline. Solver, device, thread-count and CFL
comparisons change one declared axis while retaining the case controls needed
by that experiment. Virtual p53 and p24 MCA runs form a separate stochastic
matrix and are not identified with IEEE fp64 and fp32 runs.

\section{Reference hierarchy and validation gates}
\label{sec:ch3-reference-hierarchy}

Validation claims follow the strongest available reference. Brio--Wu uses an
aligned fp64 numerical reference at $N=8000$; it is not treated as an exact MHD
solution. For the two-dimensional cases, each fine grid is conservatively
block-averaged before comparison with the adjacent coarse grid. These
self-reference differences show the direction of a three-grid trend but do not
by themselves establish an asymptotic convergence regime. Completion requires
the declared final time and expected output, while the physical-state gate
requires finite conservative variables with positive density and thermal
pressure. Transverse invariance, CPU/GPU saved-state equality, and divergence
control are separate implementation gates rather than substitutes for a
solution reference.

The Orszag--Tang configuration follows \citet{toth2000divb} after the coordinate
change $X=2\pi x$, $Y=2\pi y$ from $[0,2\pi]^2$ to the periodic unit square. The
corresponding time is $T=2\pi t$, so the reported project time $t=0.5$ maps to
$T=\pi$, the late comparison time used in that source to rounding precision.

\section{Metrics and statistical treatment}
\label{sec:ch3-metrics}

For a same-grid difference $d_{ij}=q_{ij}-q^{\mathrm{ref}}_{ij}$ over
$N=N_xN_y$ cells, the reported mean norms are
\begin{equation}
 L_{1,\mathrm{mean}}=\frac{1}{N}\sum_{ij}|d_{ij}|,\qquad
 L_{2,\mathrm{mean}}=\left(\frac{1}{N}\sum_{ij}d_{ij}^{2}\right)^{1/2},
 \label{eq:ch3-mean-norms}
\end{equation}
with $L_\infty=\max_{ij}|d_{ij}|$. They are distinct from physical-domain
integrals weighted by $\Delta x\Delta y$.
For the compact cross-system matrix, the mean-relative density difference is
\begin{equation}
 L_{1,\mathrm{rel}}(\rho)=
 \frac{N^{-1}\sum_{ij}|\rho_{ij}-\rho^{\mathrm{ref}}_{ij}|}
      {N^{-1}\sum_{ij}|\rho^{\mathrm{ref}}_{ij}|},
 \label{eq:ch3-mean-relative}
\end{equation}
and is used only when the denominator is finite and positive. This
normalisation permits scale-aware reporting within the compact matrix, but it
does not turn an fp64 baseline into an exact solution or support a ranking of
different physical systems.
Figure~\ref{fig:ch4-resolution-precision} uses absolute density
$L_{1,\mathrm{mean}}$ after block averaging. Its engineering diagnostic is
\begin{equation}
 p_{\mathrm{obs}}=\log_2\!\left(\frac{E_{128,256}}{E_{256,512}}\right).
 \label{eq:ch3-observed-p}
\end{equation}
For the Chapter~5 precision context, let
$D_N=\|\rho_{32,N}-\rho_{64,N}\|_{1,\mathrm{mean}}$ and let
$E^{64}_{256,512}$ be the matched fp64 adjacent-grid mean-$L_1$ difference
after block averaging. The reported ratio
$S_N=D_N/E^{64}_{256,512}$ compares precision discrepancy with the finest
available refinement scale within one case and solver. It is not an exact-error
ratio or an fp32-adequacy criterion.

The divergence diagnostic applies centred differences on interior cells,
\begin{align}
 D_{ij} &=\frac{B_{x,i+1,j}-B_{x,i-1,j}}{2\Delta x}
          +\frac{B_{y,i,j+1}-B_{y,i,j-1}}{2\Delta y}, \notag\\
 \mathrm{divB}_{\mathrm{mean}}
        &=\frac{1}{N_I}\sum_{(i,j)\in I}|D_{ij}|,
 &\mathrm{divB}_{\max}
        &=\max_{(i,j)\in I}|D_{ij}|.
 \label{eq:ch3-divb}
\end{align}
where $I$ excludes the outer stencil layer. These are raw, grid-scale
quantities rather than dimensionless divergence errors; mean and maximum values
are therefore interpreted separately. Same-precision CPU/GPU comparisons use
absolute $L_\infty$ and maximum ULP distance on the saved conservative state.
For arrays with the same IEEE type, sign-adjusted bit patterns are mapped to
monotone integers and the largest cellwise integer distance is reported. Thus
zero ULP denotes bitwise equality; ULP distance is not used between fp32 and
fp64 arrays and is not a physical error norm.

For $n$ MCA samples, the per-cell sample standard deviation is
\begin{equation}
 \sigma_j=\left[\frac{1}{n-1}\sum_{s=1}^{n}
 (q_{s,j}-\bar q_j)^2\right]^{1/2},\qquad
 \mathrm{spread}_q=\max_j\sigma_j .
 \label{eq:ch3-mca-spread}
\end{equation}
The maximum is a conservative spatial aggregation for shocks and shear layers.
The associated numerical signal-to-noise ratio is
$\operatorname{mean}_j|\bar q_j|/\operatorname{mean}_j\sigma_j$; it is not a
physical turbulence SNR. Sample count, virtual precision, solver, grid and
final time accompany every MCA value. Results from unmatched scopes are not
ratioed or merged.

Elapsed time is subprocess wall time from solver launch through completion and
required binary output, rather than kernel time. Repeated timing groups discard
one warm-up and retain five measurements. They report the median and
$\mathrm{IQR}=Q_{75}-Q_{25}$; speed-up is the matched baseline median divided
by the comparison median. These measurements support only the recorded
workstation and configuration.

The temporal experiment aligns fp32 and fp64 saved states at each time and
computes density $L_{1,\mathrm{mean}}$ and $L_\infty$. For positive samples in
the predeclared case-specific window, ordinary least squares fits
\begin{equation}
 \log e(t)=a+\lambda t .
 \label{eq:ch3-temporal-fit}
\end{equation}
Fit quality is reported using $R^2$, root-mean-square log residual and maximum
absolute log residual without changing the fixed window. The slope is a
bounded Lyapunov-like engineering diagnostic, not a formal maximal Lyapunov
exponent or physical instability rate.

\section{Harness, metadata, and retention}

Every evidence packet follows \texttt{config -> build -> run -> measure ->
aggregate -> plot}. Generated configurations and run metadata record the
executable and configuration hashes, process return code, declared final time,
diagnostics and elapsed wall time. Aggregators check matrix completeness,
finite admissible states and experiment-specific gates before writing machine-
readable JSON/CSV and report-facing figures. Logs and summaries are retained;
large transient grids are removed except where a retained grid is required to
close and revalidate a declared comparison pair.

\section{Deliberate exclusions}

MPI is excluded because the evaluated solver path and evidence matrix do not
provide a matched MPI implementation whose ordering effects could be isolated.
HLLD-on-GPU, KH-on-GPU, full-grid KH MCA and a second timing machine are also
outside the completed matrix. Their omission limits the claims; it does not
imply that parallel ordering, hardware or solver choice is generally
irrelevant.
